\documentclass[11pt,twocolumn]{article}
\usepackage[margin=1in]{geometry}
\usepackage{titlesec}
\usepackage{booktabs}
\usepackage{array}
\usepackage{parskip}
\usepackage{enumitem}
\usepackage{microtype}
\usepackage{times}
\usepackage[hidelinks]{hyperref}

\hypersetup{colorlinks=true, linkcolor=blue, urlcolor=blue, citecolor=blue}

\titleformat{\section}{\normalfont\large\bfseries}{}{0em}{}[\titlerule]
\titleformat{\subsection}{\normalfont\normalsize\bfseries}{}{0em}{}

\titlespacing*{\section}{0pt}{10pt}{4pt}
\titlespacing*{\subsection}{0pt}{6pt}{2pt}

\title{
    \vspace{-1.5em}
    {\large \textbf{Concept Paper}} \\[0.3em]
    {\LARGE \textbf{Quantum Circuit Classification as a}} \\
    {\LARGE \textbf{Hardware Benchmarking Framework}} \\[0.3em]
    {\normalsize for Joint Force Decision Support Applications} \\[0.5em]
    \rule{\linewidth}{0.4pt}
}

\author{
    Andrew Maciejunes \quad Ross Gore \quad Sachin Shetty \quad Barry Ezell \\
    {\small Virginia Modeling, Analysis, and Simulation Center (VMASC)} \\
    {\small Center for Secure and Intelligent Critical Systems (CSICS)} \\
    {\small Old Dominion University, Norfolk, VA}
}

\date{}

\begin{document}

\maketitle
\thispagestyle{empty}
\vspace{-1em}

% -------------------------------------------------------
\section{Problem Statement}
% -------------------------------------------------------

As the Department of Defense evaluates quantum computing
for use in wargaming, simulation, and decision support
applications, a fundamental verification problem emerges:
\textit{how does a commander or analyst know whether a
quantum device is actually performing non-classical
computation, rather than behaving as an expensive
classical simulator?}

Existing benchmarking methods address gate-level noise
(randomized benchmarking), aggregate scalar metrics
(quantum volume), and entanglement verification
(entanglement witnesses), but none directly answers the
operational question: \textit{is this device operating in
the correct computational complexity class for the task
assigned to it?} A device with low gate error rates can
still produce classically simulable outputs due to
systematic miscalibration. A device that passes an
entanglement witness can still fail to execute non-classical
algorithms correctly. Neither method is sufficient for
mission-critical validation.

% -------------------------------------------------------
\section{Proposed Approach}
% -------------------------------------------------------

We propose a \textbf{distribution-level quantum hardware
benchmark} grounded in a peer-reviewed foundation
established in our prior work~\cite{maciejunes2026}. The
core insight is that quantum circuit families occupy
distinct computational complexity classes, each producing
statistically distinguishable output distributions:

\begin{itemize}[leftmargin=*, itemsep=2pt]
    \item \textbf{Clifford circuits} -- classically
    simulable via the stabilizer formalism \cite{gottesman1997stabilizer}; output
    distributions are fully characterizable by classical
    means.
    \item \textbf{Clifford+T circuits} -- transitional
    regime; T-gate density breaks classical simulability \cite{boykin1999universal},
    producing distributions intermediate in complexity.
    \item \textbf{IQP circuits} -- believed classically
    hard to sample from \cite{bremner2016average}; output distributions are
    concentrated in local, nearest-neighbor correlations
    in the computational basis.
\end{itemize}

Our prior work demonstrated that these three families
can be classified using only polynomial measurement
resources---specifically, Z-basis pairwise correlations
extracted from $s = 16n^2$ measurement shots---achieving
a peak Random Forest accuracy of $0.91$ at 4--5 qubits.
A formal theoretical framework establishes that this
result is a provable consequence of the algebraic
structure of IQP circuits, not an artifact of the
experimental setup.

\subsection{Hardware Validation Protocol}

The proposed benchmarking protocol proceeds as follows.
Given a target quantum device, we execute a suite of
IQP, Clifford, and Clifford+T circuits of varying depth.
We collect measurement shots from the device and extract
the same polynomial feature set used in our classifier:
Hamming-weight histograms, single-qubit marginals, and
pairwise Z-basis correlations. Our pre-trained classifier
then assigns each circuit's output distribution to one of
the three complexity classes.

\textbf{Interpretation:} If a device running IQP circuits
is classified as Clifford-like, its output distribution
has degraded to the classically simulable regime---the
hardware is not performing the intended non-classical
computation. If the classifier correctly identifies the
circuit family, the device is producing outputs
statistically consistent with the target complexity class.
This provides a \textit{task-relevant, distribution-level
certificate of computational capability} that complements
existing gate-level diagnostics.

\subsection{Integration with Existing Benchmarking Methods}

The proposed method is designed to be \textbf{complementary},
not competitive, with established protocols. Table~\ref{tab:comparison}
summarizes how each method contributes to a complete
benchmarking picture.

\begin{table}[h!]
\centering
\small
\setlength{\tabcolsep}{4pt}
\caption{Benchmarking method comparison.}
\label{tab:comparison}
\begin{tabular}{p{2.0cm} p{2.5cm} p{2.0cm}}
\toprule
\textbf{Method} & \textbf{What It Certifies} & \textbf{Limitation} \\
\midrule
Randomized Benchmarking \cite{knill2008randomized} & Gate error rate & Not task-relevant \\ \hline
Quantum Volume \cite{blume2020volumetric} & System-wide scalar & Single metric  \\ \hline
Entanglement Witnesses \cite{amaro2020design} & Entanglement present & Only for select entangled states \\ \hline
Cross-Entropy (XEB) \cite{aaronson2019classical} & Output fidelity & Exponential classical cost \\ \hline
\textbf{Our Method} \cite{bremner2016average} & \textbf{Complexity class of output} & \textbf{Degrades above $\sim$12 qubits} \\
\bottomrule
\end{tabular}
\end{table}

A complete hardware assessment pipeline would first apply
randomized benchmarking to identify miscalibrated gates,
then quantum volume to establish a system-wide capability
scalar, and finally our distribution classifier to certify
that the device is producing outputs in the intended
computational complexity class. Each method diagnoses a
different failure mode; together they provide defense-grade
confidence in hardware readiness.

\subsection{Deep Learning Extension}

Our prior work established that tree-based classifiers
(Random Forest, Decision Tree) outperform linear methods
(Logistic Regression, SVM) on this task, confirming that
the class boundaries are nonlinear in the feature space.
The proposed effort will extend the classification pipeline
to include deep learning architectures such as
multilayer perceptrons and 1D convolutional networks
operating on the raw correlator feature vectors. This
extension serves two purposes: (1)~testing whether deep learning recovers
discriminative structure at qubit counts where classical
classifiers degrade, which is the primary open question
identified in our prior work, and (2)~establishing a well rounded classification pipeline that can be deployed on
operational hardware with the best possible accuracy.

% -------------------------------------------------------
\section{Current Limitations}
\label{sec:limitations}
% -------------------------------------------------------

We present these limitations transparently, as each
defines a concrete research direction for the proposed
effort.

\textbf{Scaling collapse at $\sim$12 qubits.} All four
measurement strategies degrade to near-chance accuracy
($\approx 0.33$) above approximately 12 qubits under the
quadratic shot budget $s = 16n^2$. Operational quantum
hardware of current interest (IBM, IonQ, Quantinuum)
operates at 50--400+ qubits. The method as currently
implemented does not scale to these regimes.

\textbf{Known cause: The collapse is not a fundamental
impossibility result.} It arises because the quadratic
shot budget provides a constant $\approx 32$ shots per
correlator regardless of qubit count, while the
distributional differences between circuit families
become increasingly subtle as system size grows. The
knob is known; it has not yet been turned.

\textbf{Noiseless simulation only.} All prior results
use noiseless classical simulation (Qiskit Aer). Real
device noise will alter output distributions in ways
that may be family-dependent, potentially complicating
classification. This is the most immediate gap between
current results and operational applicability.

\textbf{Linear non-separability.} The three circuit
families are not linearly separable in the Pauli
correlator feature space, as evidenced by near-chance
logistic regression performance. Classification requires
nonlinear methods, which adds inference complexity for
resource-constrained operational deployment.

% -------------------------------------------------------
\section{Research Roadmap and Funding Priorities}
% -------------------------------------------------------

The following research thrusts directly address the
limitations above and constitute the proposed effort.

\textbf{Thrust 1: Superquadratic shot budget scaling.}
Increasing the shot budget exponent from $s = \lambda n^2$
to $s = \lambda n^3$ or $s = \lambda n^4$ would provide
more shots per correlator as system size grows,
potentially shifting the collapse threshold to
operationally relevant qubit counts. The computational
cost of additional shots scales polynomially and is
well within reach of current cloud quantum hardware
access models.

\textbf{Thrust 2: Local feature selection.} Our prior
work showed that nearest-neighbor ZZ correlations
(O$(n)$ features) match all-pairs performance to within
0.02 in Random Forest accuracy. A feature set that grows
linearly rather than quadratically concentrates the shot
budget on the most discriminative correlators.
Systematically characterizing which feature subsets
preserve discriminative power at large qubit counts is
a tractable and high-payoff research direction.

\textbf{Thrust 3: Deep learning on raw measurement data.}
Classical classifiers operate on hand-engineered
correlator features. A neural network operating directly
on raw bitstring samples or shot-level data may identify
distributional structure that survives to larger qubit
counts without requiring explicit feature engineering.
This is the primary motivation for the deep learning
component and will be the first empirical test of
whether the scaling collapse is a feature extraction
problem or a fundamental information-theoretic limit.

\textbf{Thrust 4: Noisy hardware experiments.} Execution
on real quantum hardware (IBM Quantum, IonQ, or
Quantinuum via existing DoD access agreements) will
establish ground truth for how device noise affects
classification accuracy across circuit families. This
is the direct path from proof-of-concept to operational
validation and is the highest-priority experimental
next step.

% -------------------------------------------------------
\section{Summary}
% -------------------------------------------------------

We have established, through peer-reviewed work, that
quantum circuit families can be classified using
polynomial measurement resources, with a formal
theoretical guarantee that the result is not an
experimental artifact. We propose to extend this
foundation into a hardware benchmarking framework that
directly certifies whether a quantum device is operating
in the intended computational complexity class---a
capability that no existing benchmarking method provides.
Current limitations are well-understood, their causes
are identified, and the proposed research thrusts
address each one concretely. This effort positions
ODU/VMASC as a foundational contributor to DoD quantum
hardware validation methodology, with direct relevance
to J7's mandate to advance the operational effectiveness
of the future joint force through rigorous capability
assessment.

\vspace{0.5em}
\noindent\rule{\linewidth}{0.4pt}
{\small
\textbf{Point of Contact:} Andrew Maciejunes,
\texttt{amacieju@odu.edu} \\
\textbf{Institution:} Old Dominion University,
Virginia Modeling, Analysis, and Simulation Center \\
\textbf{Classification:} UNCLASSIFIED // FOR OFFICIAL USE
}


\bibliographystyle{IEEEtran}
\bibliography{references}
\end{document}